\subsection{Irreducible Representations}

In this issue, G is a topological group.

DEFINITION 2. --- A representation \(\varrho\) of G in a topological K-vector space E is said to be irreducible if \(\{0\}\) is not dense in E and if the only subrepresentations of \(\varrho\) are \(\varrho\) and the representation in the closure of \(\{0\}\) .

If \(\varrho\) is an irreducible representation of G in a separated topological vector space E, then every nonzero element of E is a cyclic vector for \(\varrho\) .

Lemma 2. --- Let \(\pi\) and \(\varrho\) be continuous linear representations of G in separated topological K-vector spaces. We assume that \(\pi\) is irreducible. Every nonzero G-morphism from \(\pi\) to \(\varrho\) is injective and every nonzero G-morphism from \(\varrho\) to \(\pi\) has dense image.

In particular, if \(\pi\) and \(\varrho\) are irreducible and of finite dimension, then every nonzero G-morphism from \(\pi\) to \(\varrho\) is an isomorphism.

Since the space of \(\varrho\) is separated, the kernel of a G-morphism \(u\) from \(\pi\) to \(\varrho\) is closed, and hence defines a subrepresentation of \(\pi\) . If the morphism \(u\) is not zero, its kernel must therefore be reduced to \(\{0\}\) , since \(\pi\) is an irreducible representation in a separated topological vector space. Similarly, the closure of the image of a nonzero G-morphism from \(\varrho\) to \(\pi\) is a nonzero subrepresentation of \(\pi\) , hence is equal to the space of \(\pi\) .

The last assertion follows from what precedes.

Lemma 3. --- Let \(\varrho\) be a continuous linear representation of G in a separated topological K-vector space E of finite dimension. If E is nonzero, then there exists an irreducible subrepresentation of \(\varrho\) .

Since E is of finite dimension and nonzero, there exists a nonzero G-invariant subspace F of E of minimal dimension. This subspace is closed in E and defines a subrepresentation of E; every subrepresentation of F is also a subrepresentation of E, and the representation F is therefore irreducible by minimality.

Remark. --- A nonzero representation E of G does not always contain an irreducible subrepresentation (cf. V, p. 426, remark). One will see, however, that this is the case if G is compact, if K = C, and if E is a separated quasi-complete locally convex space over K and nonzero (prop. 7 of V, p. 464).

